% ============================================
% DOC-18
% Gerado de molambudos.md — build_miolo.py
% ============================================
\section*{\textbf{DOC-18} --- Relatório Oficial do Curral}
\markboth{DOC-18 --- Relatório Oficial do Curral}{}

\noindent \textit{1916 --- Reconstituição sobre o Campo do Alagadiço}

\noindent \textbf{Reconstituição ficcional --- sem equivalente localizado no Arquivo Nacional}
\textbf{1916 --- Documento narrativo, não transcrição}

\vspace{0.5em}\noindent\rule{\textwidth}{0.5pt}\vspace{0.5em}

\noindent \textit{RECONSTITUIÇÃO FICCIONAL DE EXPEDIENTE ADMINISTRATIVO}
\textit{SOCORROS PÚBLICOS}
\textit{RELATÓRIO NARRATIVO DOS CAMPOS DO CEARÁ}
\textit{EXERCÍCIO DE 1915-1916}

\vspace{0.5em}\noindent\rule{\textwidth}{0.5pt}\vspace{0.5em}

\noindent \textbf{Campo:} Alagadiço, próximo a Fortaleza
\textbf{Capacidade oficial:} não indicada nesta reconstituição
\textbf{Número de internados durante o período:} impreciso; as fontes consultadas registram milhares de confinados
\textbf{Óbitos:} sem total verificável para reprodução nesta narrativa

\noindent \textbf{Voz do relator ficcional:}

\noindent \textit{\textquotedblleft{}Os campos de concentração, oficialmente chamados 'currais do governo', foram estabelecidos como medida emergencial para conter o fluxo de retirantes em direção à capital. Na prática, funcionaram como depósitos de seres humanos em estado de inanição.}

\noindent \textit{A alimentação e as condições sanitárias eram insuficientes para proteger a vida dos confinados. Esta reconstituição não atribui percentual de mortalidade nem autoria documental a essa formulação.}

\noindent \textit{Recomenda-se o fechamento dos campos e a adoção de medidas efetivas de socorro.\textquotedblright{}}

\vspace{0.5em}\noindent\rule{\textwidth}{0.5pt}\vspace{0.5em}

\noindent \textbf{Nota editorial da reconstituição:}

\noindent A existência do Campo do Alagadiço, próximo a Fortaleza, durante a seca de 1915 é documentada. Os números, a voz do relator e a forma administrativa acima são recursos ficcionais. Este texto não afirma que o campo permaneceu aberto até 1917; a transferência de Joaquim nesse ano pertence à cronologia ficcional do personagem.

\vspace{0.5em}
\noindent\textit{\textbf{↪ Rotas:} \rota{DOC-19} \rotasep \rota{MEM-06}}